% 02_background.tex
\section{Background}
\label{sec:background}

\subsection{Path graphs and scalar filtrations}

Let $G_n = (V, E)$ denote the \emph{path graph} on $n$ vertices, where
$V = \{1,\ldots,n\}$ represents the residues of a polypeptide chain and
$E = \{(i,\,i{+}1) : i = 1,\ldots,n{-}1\}$ the peptide bonds.
A \emph{scalar filtration} of $G_n$ assigns a real value $u_i$ to each
vertex and grows a sequence of nested subgraphs
$\emptyset = \mathcal{G}_{-\infty} \subseteq \mathcal{G}_{u_{(1)}} \subseteq
\cdots \subseteq \mathcal{G}_{u_{(n)}} = G_n$
by adding vertices in increasing $u$ order together with all incident
edges whose other endpoint is already present~\cite{edelsbrunner2010computational}.

Following~\cite{cazals2025alphafold}, we set $u_i = -s_i$ where $s_i \in [0,100]$
is the pLDDT score of residue $i$.
Processing by increasing $u$ therefore inserts residues in
\emph{decreasing pLDDT order}: high-confidence residues enter first.
An edge $(i,i{+}1)$ is activated as soon as both its endpoints have been
inserted.
We denote this filtration $\mathcal{G}_{\mathrm{pLDDT}}$.

\subsection{Connected components and the Elder Rule}

Let $N_{cc}(u)$ count the connected components of $\mathcal{G}_u$ at
parameter $u$.
Initially $N_{cc} = 0$; each vertex insertion increments it by 1 (a new
singleton component is born), and each edge activation that joins two distinct
components decrements it by 1 (one component dies by merging into the other).

The \textbf{Elder Rule} determines which component survives a merge:
the component with the \emph{lower} $u$-value at birth (equivalently, born at
\emph{higher} pLDDT) is the elder and persists; the younger component dies,
recording a \emph{persistence pair} $(b, d)$ where $b = \mathrm{birth\_pLDDT}$
of the younger component and $d = \mathrm{death\_pLDDT}$ (the pLDDT level at
which the merge occurs).
Persistence in pLDDT units is $p = b - d \geq 0$.

The collection of all $n{-}1$ persistence pairs constitutes the
\textbf{persistence diagram} (PD)
$\mathcal{P} = \{(b_i, d_i)\}_{i=1}^{n-1}$~\cite{edelsbrunner2010computational}.
Pairs with $b_i = d_i$ (zero persistence) are called \emph{accretions}:
they occur when two components that were already in the same root's tree
merge, or when a batch of residues at identical pLDDT level are inserted
simultaneously and their internal edges are activated at birth level $L$
(so birth equals death equals $L$).

\subsection{Five statistics}
\label{sec:bg_stats}

We reproduce Definition~1 of~\cite{cazals2025alphafold}.
Let $m = n{-}1$ be the total number of PD pairs, $m'$ the number of pairs
with $b_i > d_i$ (strictly positive persistence), and
$\mathcal{P}_D$ the set of \emph{distinct} persistence values
$\{p_i = b_i - d_i\}$.
Write $\mathbb{P}[v] = \#\{i : p_i = v\}/m$ for the empirical probability of
value $v$.

\begin{enumerate}
  \item \textbf{Fraction of positive critical points}:
        $f^+_{cp} = m'/(n{-}1)$.
        High values indicate a heterogeneous pLDDT profile along the sequence.

  \item \textbf{Mean persistence}:
        $\bar{p} = \frac{1}{m}\sum_{i=1}^{m}(b_i - d_i)$.

  \item \textbf{Normalised persistence entropy}:
        \begin{equation}
          H_p = -\sum_{v \in \mathcal{P}_D} \mathbb{P}[v]\,\log\mathbb{P}[v]
                \;\Big/\;\log|\mathcal{P}_D|.
          \label{eq:Hp}
        \end{equation}
        $H_p \in [0,1]$; it equals $0$ when all pairs share a single
        persistence value and $1$ when all values are equally probable.
        Ordered proteins with narrow pLDDT distributions yield low $H_p$;
        multi-domain proteins with wide distributions yield high $H_p$.

  \item \textbf{Peak connected-component count}:
        $N_{cc}^{max} = \max_u N_{cc}(u)$.
        For random pLDDT the expected peak is $\approx n/4$~\cite{cazals2025alphafold}.

  \item \textbf{Persistent local maxima}:
        $\mathrm{PLM}(t_p)$ counts local maxima of $N_{cc}(u)$ whose
        persistence under the super-level-set filtration of $N_{cc}$
        (viewed as a 1D scalar function of the filtration index) exceeds
        the absolute threshold $t_\nu = n \cdot t_p$.
        A high PLM count is the topological signature of multi-domain
        proteins, where each structural domain produces a distinct
        connected-component peak.
\end{enumerate}

$f^+_{cp}$ and $H_p$ are strongly linearly correlated across any proteome
(Pearson $r \approx 0.97$ in~\cite{cazals2025alphafold}), providing a
cross-check on both statistics.

\subsection{Arity signature}
\label{sec:bg_arity}

The \emph{arity} of residue $k$ is the number of $\mathrm{C}_\alpha$ atoms
within a Euclidean distance $r = 10$\,\AA{} of $\mathrm{C}_\alpha(k)$,
excluding residue $k$ itself~\cite{cazals2025alphafold}.
This radius accounts for non-covalent contacts and has been shown sufficient
to infer protein domains via spectral clustering.

\begin{definition}[Arity signature {\cite[Def.\,2]{cazals2025alphafold}}]
\label{def:arity_sig}
Let $L_a = (a_1,\ldots,a_n)$ be the per-residue arity values.
Given quantiles $q_1 = 0.25$ and $q_2 = 0.75$, the \emph{arity signature}
of a protein $P$ is the pair
\[
  \mathrm{Sig}(P) = (a_{25},\;a_{75})
  \quad\text{where}\quad
  a_k = \mathrm{CDF}_{L_a}^{-1}(q_k).
\]
\end{definition}

The \emph{arity map} (Definition~3 of~\cite{cazals2025alphafold}) is a
2D histogram of all proteome proteins in the $(a_{25}, a_{75})$ plane.
Compactly folded multi-domain proteins cluster in high-arity bins
($a_{25} \geq 7$, $a_{75} \geq 13$); intrinsically disordered proteins
occupy low-arity bins ($a_{25} \leq 4$).
Intuitively, the arity signature compresses the full $\mathrm{C}_\alpha$
packing-density distribution into two robust quantile descriptors.
